\documentclass[10pt,letterpaper]{article}

\usepackage[T1]{fontenc}
\usepackage[utf8]{inputenc}
\usepackage{lmodern}
\usepackage[margin=0.82in,headheight=14pt]{geometry}
\usepackage{microtype,amsmath,amssymb,array,tabularx,longtable,booktabs,multicol}
\usepackage{graphicx,xcolor,enumitem,fancyhdr,hyperref,titlesec,needspace}
\definecolor{Navy}{HTML}{122A42}
\definecolor{Ocean}{HTML}{147D91}
\definecolor{Soft}{HTML}{E9F2F5}
\definecolor{Ink}{HTML}{263745}
\definecolor{Gray}{HTML}{5F6873}
\hypersetup{colorlinks=true,linkcolor=Ocean,urlcolor=Ocean,bookmarksnumbered=true}
\setlength{\parindent}{0pt}
\emergencystretch=2em
\setlength{\parskip}{5.3pt}
\setlist{topsep=3pt,itemsep=2pt,parsep=0pt,leftmargin=*}
\titleformat{\section}{\large\bfseries\color{Navy}}{\thesection}{0.65em}{}
\titleformat{\subsection}{\normalsize\bfseries\color{Ocean}}{\thesubsection}{0.6em}{}
\titleformat{\subsubsection}{\normalsize\bfseries\color{Navy}}{\thesubsubsection}{0.5em}{}
\pagestyle{fancy}\fancyhf{}
\fancyhead[L]{\textsf{Architecture Documentation \& Engineering}}
\fancyhead[R]{\textsf{\DocCategory}}
\fancyfoot[L]{\footnotesize\textcolor{Gray}{Controlled learning / engineering artifact}}
\fancyfoot[R]{\footnotesize\thepage}
\renewcommand{\headrulewidth}{0.35pt}
\newcommand{\DocCategory}{Suite}
\newcommand{\ArchTitle}[4]{%
\renewcommand{\DocCategory}{#3}%
{\small\sffamily\bfseries\color{Ocean} #3 \hfill #4}\par\vspace{7pt}
{\Large\bfseries\color{Navy}#1}\par\vspace{4pt}
{\normalsize\color{Gray}#2}\par\vspace{11pt}
\textcolor{Ocean}{\rule{\linewidth}{1.15pt}}\vspace{4pt}
}
\newcommand{\Note}[1]{\par\smallskip\noindent\fcolorbox{Ocean}{Soft}{\parbox{0.94\linewidth}{\small\textbf{Engineering note.}~#1}}\par\smallskip}
\newcommand{\Outcome}[1]{\textbf{Outcome.} #1\par}
\newcommand{\Evidence}[1]{\textbf{Required evidence.} #1\par}
\newcommand{\BlankLines}[1]{\par\begingroup\setlength{\parskip}{9pt}\count255=0\loop\ifnum\count255<#1\noindent\rule{\linewidth}{0.35pt}\par\advance\count255 by 1\repeat\endgroup\vspace{4pt}}
\newcommand{\BlankBox}[1]{\noindent\fbox{\parbox[t][#1][t]{0.98\linewidth}{\mbox{}}}\par\smallskip}
\newcommand{\Field}[2]{\Needspace{3\baselineskip}\textbf{#1}\par\textit{\color{Gray}#2}\par\BlankLines{2}}
\newcommand{\FieldBlank}[1]{\Needspace{3\baselineskip}\textbf{#1}\par\BlankLines{2}}
\newcommand{\DocProvenance}{\textit{Source foundation: Clements et al., \emph{Documenting Software Architectures: Views and Beyond}, 2nd ed. (2011). Interpretive examples, exercises, checklists, and templates are original adaptations, not reproductions.}\par}
\newcommand{\SmallHeader}[1]{\smallskip\textbf{\color{Navy}#1}\par}
\newcolumntype{Y}{>{\raggedright\arraybackslash}X}
\newcommand{\Tick}{\(\square\)\,}

\begin{document}
\ArchTitle{Prerequisite-Based Learning Curriculum}{Twelve-week progression from architectural structures to reviewed documentation}{Curriculum}{ARCH-CUR-001}
\DocProvenance
\section{Entry competencies and prerequisites}
\small
\begin{tabularx}{\linewidth}{@{}YYY@{}}
\toprule
\textbf{Prerequisite} & \textbf{Observable checkpoint} & \textbf{If not met} \\
\midrule
Programming and modularity & can name responsibilities and dependencies in a codebase & review cohesion/coupling and package structure \\
Distributed systems basics & can explain call, event, failure and consistency tradeoffs & study APIs, concurrency and messaging \\
Networking and deployment & can diagram network reachability and runtime placement & review network boundaries and cloud primitives \\
Requirements and testing & can write acceptance criteria and a testable quality scenario & practice scenario writing and validation \\
Version control and tooling & can version documentation and reproduce a build & learn Git and basic LaTeX/diagram-as-code \\
\bottomrule
\end{tabularx}
\normalsize
\section{Learning dependency graph}
Foundation -> module modeling -> runtime modeling -> allocation -> interfaces and behavior -> viewpoint specification -> multi-view description -> traceability and review. Chapter 9 view selection should be practiced early for orientation and revisited after studying the three structural families.

\section{Twelve-week staged curriculum}
\small
\begin{tabularx}{\linewidth}{@{}YYYY@{}}
\toprule
\textbf{Week} & \textbf{Primary study} & \textbf{Lab focus} & \textbf{Checkpoint} \\
\midrule
1 & Prologue; Ch. 9 preview & charter, stakeholder interviews & stakeholder/concern matrix \\
2 & Ch. 1 & module identification, catalog & module diagram \\
3 & Ch. 2 & uses dependencies and DSM & cycle analysis \\
4 & Ch. 3 & runtime components and connectors & runtime inventory \\
5 & Ch. 4 & communication pattern alternatives & connector contract \\
6 & Ch. 5 & runtime-to-infrastructure allocation & deployment mapping \\
7 & Ch. 6 & context, variability, decisions & decision record \\
8 & Ch. 7 & interface syntax and semantics & interface contract \\
9 & Ch. 8 & interaction traces and state models & failure behavior model \\
10 & Ch. 9 and Ch. 10 & view selection and packaging & architecture description v1 \\
11 & Ch. 11, Epilogue & stakeholder review and methods crosswalk & review findings and dispositions \\
12 & Appendices; capstone & tool-driven revision and evidence checks & revised AD and defense \\
\bottomrule
\end{tabularx}
\normalsize
\section{Mastery gates}
\begin{itemize}
\item Gate A after week 3: demonstrate typed module relations, change-impact reasoning, and a dependency matrix.
\item Gate B after week 6: maintain consistent IDs and mappings between module, C\&C and allocation representations.
\item Gate C after week 9: specify at least two failure paths, a provided/required interface, and a measurable quality scenario.
\item Gate D after week 12: pass an evidence-based documentation review with critical findings resolved.
\end{itemize}
\section{Grading and capstone requirements}
\small
\begin{tabularx}{\linewidth}{@{}YYY@{}}
\toprule
\textbf{Component} & \textbf{Weight} & \textbf{Minimum bar} \\
\midrule
Engineering notebooks and diagrams & 30\% & each major view family represented \\
Contracts and decision rationale & 20\% & at least two decision trade-offs \\
Traceability and consistency & 20\% & no unresolved orphan critical concern \\
Review and revision & 20\% & documented findings and follow-up \\
Technical defense & 10\% & answer three stakeholder use cases \\
\bottomrule
\end{tabularx}
\normalsize
\section{Evidence log}
For each week record date, source chapter, artifact IDs, a demonstration of understanding, mistakes found, revision commit, and next dependency. Use the traceability workbook and assessment templates for repeatable evidence.

\BlankLines{5}

\end{document}
